\section*{Chapter \ref{ch:g1:floating-and-sinking} --- Floating and Sinking}

\begin{solution}{exo:g1:floating-and-sinking:1}
Floaters: the cork, the dry twig, the plastic duck. Sinkers: the key,
the marble, the stone.
\end{solution}

\begin{solution}{exo:g1:floating-and-sinking:2}
A splash can push even a floater underwater for a moment, so a thrown
object might fool you. Putting it down gently and waiting makes the
test fair.
\end{solution}

\begin{solution}{exo:g1:floating-and-sinking:3}
The metal spoon sank while the plastic spoon floated. Nora should
look at the \emph{stuff} a thing is made of, not its name or its job.
\end{solution}

\begin{solution}{exo:g1:floating-and-sinking:4}
Heavy and floating: a tree trunk (or a ship). \omterm{def:g1:light-and-shadows:source}{Light} and sinking: a
coin (or a key). So ``heavy things \omterm{def:g1:floating-and-sinking:words}{sink}'' is a bad rule --- weight
alone does not decide.
\end{solution}

\begin{solution}{exo:g1:floating-and-sinking:5}
Yes --- wood is a floating stuff, and a thing usually behaves like
its stuff, twig or trunk. The log of
\cref{ex:g1:floating-and-sinking:log} \omterm{def:g1:floating-and-sinking:words}{floats} down the river.
\end{solution}

\begin{solution}{exo:g1:floating-and-sinking:6}
While pushing, you feel the water pushing the ball back up, harder
the deeper it goes. When you let go, the ball pops back up to the
top --- the water's push wins.
\end{solution}

\begin{solution}{exo:g1:floating-and-sinking:7}
The hollow-bowl shape: a wide hull with high sides gives the water
more to push up on. The plasticine boat of
\cref{ex:g1:floating-and-sinking:boat} shows it --- same stuff, but
the bowl shape \omterm{def:g1:floating-and-sinking:words}{floats}.
\end{solution}

\begin{solution}{exo:g1:floating-and-sinking:8}
The full closed bottle \omterm{def:g1:floating-and-sinking:words}{floats} low in the water (just barely, for most
bottles); the empty capped bottle \omterm{def:g1:floating-and-sinking:words}{floats} high, riding on the water
--- it is mostly air inside.
\end{solution}

\begin{solution}{exo:g1:floating-and-sinking:9}
The water never gets tired. Squashed into a ball, the plasticine
gives the water only a small round bottom to push on; shaped as a
hollow boat, the same lump spreads wide and the water's push on all
that bottom is enough to hold it up. It is the shape that changed,
not the water.
\end{solution}
